\begin{description}
\item[Tovan sums] The tovan sum of two whole numbers a and b is written a ⊕ b and equals a + b + a×b. Rules: (1) order does not matter: a ⊕ b = b ⊕ a. (2) 0 is neutral: 0 ⊕ b = b. (3) -1 swallows everything: (-1) ⊕ b = -1. (4) A chain is worked left to right. Worked example: 2 ⊕ 3 = 2 + 3 + 6 = 11.\\ \emph{Sample TEST item:} What is 1 ⊕ 4 (tovan sum)? A) 5; B) 9; C) 4; D) 10 (key B).
\item[Krell steps on a grid] A krell step moves a grid point (x, y) to (y, x + 1): swap the two numbers, then add 1 to the new second one. Rules: (1) a krell step always lands on a grid point. (2) Repeating krell steps from (0, 0) climbs a staircase (0,0) → (0,1) → (1,1) → (1,2) → (2,2). (3) Two krell steps move (x, y) to (x + 1, y + 1). Worked example: one krell step from (3, 5) gives (5, 4).\\ \emph{Sample TEST item:} Where does one krell step take (2, 7)? A) (7, 2); B) (3, 7); C) (8, 2); D) (7, 3) (key D).
\item[Vorn sequences] The vorn sequence starts 1, 3 and each next term is the previous term plus twice the term before it: v(n) = v(n-1) + 2×v(n-2). Rules: (1) every term is odd. (2) Each term is about twice the one before. Worked example: v(3) = 3 + 2×1 = 5, v(4) = 5 + 2×3 = 11.\\ \emph{Sample TEST item:} What is the 5th vorn term v(5)? A) 16; B) 21; C) 13; D) 19 (key B).
\item[Glint numbers] A whole number is glint when its last digit is not 0 and the number is divisible by its last digit. Rules: (1) every number ending in 1 is glint. (2) A number ending in 0 is never glint. (3) A number ending in 5 is always glint. Worked example: 24 is glint because 24 ÷ 4 = 6; 26 is not glint because 26 ÷ 6 is not whole.\\ \emph{Sample TEST item:} Which number is glint? A) 38; B) 36; C) 34; D) 27 (key B).
\item[Mirthful triangles] A triangle is mirthful when twice its longest side is at most three times its shortest side (longest ≤ 1.5 × shortest). Rules: (1) every equilateral triangle is mirthful. (2) Scaling a triangle does not change whether it is mirthful. (3) The middle side does not matter. Worked example: sides 4, 5, 6 are mirthful (2×6 = 12 ≤ 3×4 = 12); sides 3, 4, 5 are not (10 > 9).\\ \emph{Sample TEST item:} Which triangle is mirthful? A) sides 4, 4, 7; B) sides 3, 4, 5; C) sides 2, 3, 4; D) sides 5, 6, 7 (key D).
\item[Sorrel counts of lattice rectangles] For a rectangle drawn on grid points, its sorrel count is the number of grid points strictly inside it minus the number of grid points on its edges. A w × h rectangle has (w-1)×(h-1) points inside and 2×(w+h) on its edges. Rules: (1) small rectangles have negative sorrel counts. (2) Swapping width and height keeps the count. Worked example: a 3 × 3 square has 4 inside and 12 on the edges, sorrel count -8.\\ \emph{Sample TEST item:} What is the sorrel count of a 4 × 4 square? A) -16; B) -7; C) 7; D) 9 (key B).
\item[Quill distance] The quill distance between two grid points counts each horizontal unit as 1 and each vertical unit as 2: |x1 - x2| + 2×|y1 - y2|. Rules: (1) it is the same in both directions. (2) Moving up is twice as costly as moving sideways. (3) The points at quill distance 2 from a point are 2 sideways or 1 up or down from it. Worked example: from (0, 0) to (3, 2) the quill distance is 3 + 4 = 7.\\ \emph{Sample TEST item:} What is the quill distance from (0, 0) to (2, 3)? A) 7; B) 8; C) 5; D) 10 (key B).
\item[Dapple numbers] A whole number is dapple when its tens digit is greater than its units digit (a one-digit number has tens digit 0, so it is never dapple). Rules: (1) every multiple of 10 from 10 to 90 is dapple. (2) Reversing the two digits of a dapple number with different digits gives a number that is not dapple. Worked example: 52 is dapple (5 > 2); 25 is not (2 < 5); 44 is not (4 = 4).\\ \emph{Sample TEST item:} Which number is dapple? A) 37; B) 7; C) 77; D) 73 (key D).
\item[Brisk folds of a list] A brisk fold folds a list of numbers in half and adds the numbers that land on each other: first + last, second + second-to-last, and so on; in a list of odd length the middle number stays alone at the end. Rules: (1) a fold keeps the total. (2) A fold roughly halves the length. Worked example: brisk fold of [1, 2, 3, 4] is [5, 5]; of [1, 2, 3] it is [4, 2].\\ \emph{Sample TEST item:} What is the brisk fold of [2, 5, 1, 3, 4]? A) [6, 8]; B) [7, 4, 4]; C) [6, 8, 1]; D) [6, 9, 1] (key C).
\item[Orbit words] An orbit word grows by rewriting every letter at once: each X becomes XYY and each Y becomes YX. Start from X. Rules: (1) every orbit word starts with X. (2) Each rewrite makes the word longer. Worked example: X → XYY → XYYYXYX.\\ \emph{Sample TEST item:} What does X become after one rewrite? A) XXY; B) YX; C) XYY; D) XY (key C).
\item[Vells of 2×2 grids] The vell of a 2 × 2 grid of numbers [[a, b], [c, d]] is a×d minus b minus c. Rules: (1) only the top-left and bottom-right numbers are multiplied. (2) Swapping b and c keeps the vell. Worked example: the vell of [[2, 1], [3, 4]] is 2×4 - 1 - 3 = 4.\\ \emph{Sample TEST item:} What is the vell of [[3, 1], [2, 5]]? A) 17; B) 15; C) 13; D) 12 (key D).
\item[Lumen clocks] A lumen clock has 7 marks, 0 to 6, around a circle. One lumen tick moves the hand forward by 2 marks, wrapping past 6 back to 0. Rules: (1) after 7 ticks the hand is back where it started. (2) The hand visits every mark before it returns. Worked example: from mark 3, two ticks go 3 → 5 → 0.\\ \emph{Sample TEST item:} From mark 0, where is the hand after 1 tick? A) 0; B) 7; C) 1; D) 2 (key D).
\end{description}
